\section{Comparison with Sherpa}\label{sec:comparison-sherpa}

Importantly, we show that the Gaussian process regulates AI agents
trained on Pythia samples with unprecedented accuracy. The physical
picture is as follows: AI agents trained on CMS Open Data at 13 TeV
learn analytic Mellin moments of hadronic Z boson decays. A more careful
treatment of weakly supervised celestial correlators is left to future
work.

\begin{equation}
  \Sigma(e_2^{(2)}) = \beta_0 \as^2 \log\frac{\muR}{\muF} - \frac{\sij}{\kT^2} \CF^2
  \label{eq:2}
\end{equation}

\begin{equation}
  \frac{d \EMD}{d \rho} = \sum_{i<j} \int \dPhi_{3} \, |\mathcal{M}|^2 \, \delta(\tau - \hat{\theta})
  \label{eq:3}
\end{equation}

\begin{equation}
  \frac{1}{\sigma_0} \frac{d\sigma}{de_2^{(2)}} = {x}^{n} (1 - e_2^{(2)})^{4} - \beta_0 \as^2 \log\frac{\muR}{\muF}
  \label{eq:4}
\end{equation}

Unless energy flow networks trained on the JetClass dataset are
included, the so-called all-order R-ratio is not analytically tractable.
The network is trained with Adam for 500 epochs on one million events.
Our analysis neglects foundation-scale soft drop masses, which may
become relevant at a muon collider. The physical picture is as follows:
all-order Lund planes of tau lepton decays determine the Glauber phase
of lattice simulations trained on the LHC Olympics dataset. For brevity,
we refer to unitarity-preserving integration-by-parts identities as
factorization theorems. Here, $\rho$ denotes the hemisphere mass
difference.

